\section{Introduction}\label{sec:intro}

A bank that deploys a single credit-scoring or forecasting model into more than one market is now subject to several distinct model-risk regimes at once. In the United States, model risk management has been governed since 2011 by the joint Federal Reserve / OCC guidance SR 11-7 \citep{FedSR117, OCC2011}. In the European Union, the AI Act \citep{EUAIAct} imposes record-keeping, logging, and retention duties on high-risk AI systems, with the high-risk provisions taking effect on 2 August 2026. The United Kingdom's Prudential Regulation Authority issued its first dedicated model-risk supervisory statement, SS1/23 \citep{PRASS123}, effective 17 May 2024. Singapore's Monetary Authority published the FEAT principles \citep{MASFEAT} and the Veritas assessment methodologies \citep{MASVeritas}. Canada's OSFI finalised Guideline E-23 \citep{OSFIE23} in September 2025, effective 1 May 2027, expanding its scope to all models including machine learning. Two voluntary frameworks (the NIST AI Risk Management Framework \citep{NIST_AI_100_1} and ISO/IEC 42001 \citep{ISO42001}) sit alongside the binding regimes.

These regimes overlap heavily in intent but differ in letter. All require some form of documentation, validation, and governance; only the EU AI Act fixes a numerical record-retention floor; only OSFI E-23 makes model \emph{decommissioning} an explicit lifecycle stage; fairness assessment is central to the EU and Singapore regimes but absent from the original SR 11-7. A compliance officer reconciling these texts, and a researcher attempting to measure whether published work would survive any of them, both face the same obstacle: there is no common yardstick. The fragmentation is not academic: machine learning is now a core methodology in quantitative finance \citep{GuKellyXiu2020}, and the same model is routinely deployed across jurisdictions.

This paper supplies that yardstick and makes two paired contributions. First (\Cref{sec:crosswalk}), we construct a \emph{cross-jurisdiction crosswalk}: a single table mapping twelve audit-trail and model-risk-management requirements onto the specific provision of each of the five regulatory regimes and the two voluntary frameworks, so that the table can be read down a column for a single regime's full set of obligations or across a row to compare the regimes on one requirement. Second (\Cref{sec:benchmark}), we release an \emph{open compliance benchmark} that scores a machine-learning pipeline against the union of these requirements, with a reproducible scoring harness and a public leaderboard, so that compliance against the combined standard becomes a measurable, comparable quantity rather than a manual reading exercise.

\subsection{Contributions}\label{sec:contributions}

\begin{itemize}[leftmargin=*]
    \item A twelve-requirement crosswalk of machine-learning model-risk regulation across five regulatory regimes and two voluntary frameworks (\Cref{tab:crosswalk}), with each cell tied to the citable provision and a marking of whether each requirement is present, implied, partial, or absent in each regime.
    \item A structured analysis of cross-regime convergence and divergence: which requirements are universal, which are jurisdiction-specific, and which create genuine multi-market compliance conflicts.
    \item An open, versioned compliance benchmark and scoring harness, released under permissive licences, that scores an ML pipeline against the union standard and against any single regime, with a public leaderboard seeded on public-data pipelines.
\end{itemize}

\subsection{Scope and limitations}\label{sec:scope}

The crosswalk is an \emph{interpretive} mapping of public regulatory text, not legal advice; firms remain responsible for their own compliance determinations with their supervisors. Regulations evolve: the U.S.\ replaced SR 11-7 / OCC 2011-12 with new interagency guidance in 2026, and OSFI E-23 only takes effect in 2027. The crosswalk is therefore versioned to the texts as of the dates recorded in \Cref{sec:crosswalk}, and the released artefact is designed to be updated as the law changes.

\subsection{Roadmap}\label{sec:roadmap}

\Cref{sec:crosswalk} describes the crosswalk methodology, presents the mapping, and analyses convergence and divergence. \Cref{sec:benchmark} describes the open compliance benchmark, its scoring harness, and the seed leaderboard. \Cref{sec:discussion} discusses use, limitations, and future maintenance.
